\section{Experiment}
\label{sec:experiment}
We apply our framework to questions about U.S. and Japanese political parties on major GSE models, revealing their citation patterns and attack surfaces based on \emph{conjection injection barriers}.

\subsection{Experimental Setup}

\textbf{Tatget GSE Models.}
We employ three GSE models for answer generation: OpenAI GPT-5~\cite{chatgpt}, Claude Sonnet $4$ (claude-4-sonnet-20250514)~\cite{claude}, and Gemini Flash $2.5$ (gemini-2.5-flash)~\cite{gemini}, the most advanced publicly available models with web search mode enabled.
We clarify why we focus on closed GSE models in this study.
According to surveys~\cite{menlo-ventures2025stategenaienterprise, a16z2025stateconsumerai}, closed models dominate the LLM market in both enterprise and consumer segments, with Claude, OpenAI, and Gemini together accounting for 88\% of spending share, while open-source models account for only 11\%.
These models are developed in the U.S., but they dominate usage in Japan as well.
According to GMO Research \& AI~\cite{gmo-research2025japanai}, Japanese users predominantly rely on closed models such as GPT, Gemini, and Copilot, all of which are U.S.-developed products.
While Japan has developed numerous LLM models, to our knowledge, no Japanese GSE models are publicly available.
Furthermore, widely deployed GSEs such as Google AI Overview, which integrates Google Search by default, rely on the closed model Gemini~\cite{stein2025expandingaioverviewsaimode, reid2025aisearchbeyondinformation}; consequently, users routinely encounter generated answers from closed GSE models when conducting web searches, underscoring the central role these models play in contemporary information retrieval.
Note that our evaluation is not restricted to closed models, as it relies only on the inputs and outputs of GSEs (questions and answers), not their internal states.
For reproducibility, the question-answer datasets are available in our repository~\cite{QL_CIKM_RERO}.

In our experiments, the temperature of each GSE model cannot be set because APIs do not allow setting temperature for LLMs with search mode enabled; thus, the cited web sources and textual content of answers vary across generations.
To capture answer variation, we ask each question three times to suppress bias from a single answer.
The answers are obtained on May $5$, $2026$.


\textbf{User Profiles.}
We design three user profiles based on two user attributes: \textbf{political knowledge}, which represents the user's familiarity with political issues and takes the levels \textsc{Unknown} and \textsc{High}; and \textbf{political ideology}, which represents the user's political leanings and takes the levels \textsc{Conservative}, \textsc{Progressive}, and \textsc{Neutral}.

Using these attributes, we construct three user profile combinations: \textsc{Unknown-Neutral}, \textsc{Unknown-Progressive}, and \textsc{High-Progressive}.
The \textsc{Unknown-Neutral} profile represents a user with no political knowledge and a neutral ideology; the \textsc{Unknown-Progressive} profile represents a user with no political knowledge but a progressive ideology; and the \textsc{High-Progressive} profile represents a user with high political knowledge and a progressive ideology.

\textbf{Questions.}
We design political questions on topics common to both countries to elicit policy positions in a politically neutral manner, focusing on whether GSEs cite contents articulating positions on the target topics rather than favoring any particular party.

All questions are closed-ended and include a party name according to PoisonedRAGs study (e.g, ``Regarding government debt, does {PARTY} currently prioritize debt restraint or growth-oriented investment?'').

To avoid bias, we build the question set on MARPOR (Manifesto Project), an established political-science framework for classifying party manifestos.
From the four voter-facing MARPOR domains most relevant to everyday political judgment (External Relations, Economy, Welfare and Quality of Life, Fabric of Society), we select 11 topics on which all target parties take an explicit stance; the full list is given in Appendix~\ref{app:appendix_topics}.
Using these 11 topics as input, we prompt three closed-weight LLMs to generate candidate questions and curate the final set of 33.
Questions and answers are written in English for U.S. parties and in Japanese for Japanese parties.
The full question set is available in our repository~\cite{QL_CIKM_RERO}.

Following the DETAIL framework of Kim et al.~\cite{kim2025mattersmeasuringimpactprompt}, we fix question specificity at the least specific domain-level (asking about a stance on an entire policy domain without proper nouns or concrete figures).
This prevents questions and citations from being skewed toward topics that specific publishers actively promote, and reflects the abstraction level typical of voters' everyday political judgments.


\textbf{Target Political Parties.}
We target political parties satisfying each country's requirements for national political parties as primary information sources.
Specifically, in Japan we target eight parties. The ruling parties are ``Liberal Democratic Party (LDP)'' and ``Japan Innovation Party (JIP)'', and the opposition parties are ``Centrist Reform Coalition (CRA)'', ``Democratic Party for the People (DPP)'', ``Japanese Communist Party (JCP)'', ``Reiwa Shinsengumi (Reiwa)'', ``Sanseito'', and ``Conservative Party of Japan (CPJ)''. In the U.S., we target five parties. The ruling party is ``Republican Party (GOP)'', and the opposition parties are ``Democratic Party (DP)'', ``Green Party (GPUS)'', ``Libertarian Party (LP)'', and ``Constitution Party (COP)''.


\subsection{Experimental Results}


\textbf{Distribution of Content-injection Barriers.}

まず、ペルソナ別 + alignment別をマージしたチャートを出す。
モデル毎に差があります。
国別でも差があります。
これを1枚のFigureで伝える。
LLMの結果は全般、Web検索よりも一次情報の割合が高い傾向にある。。

次に野党と与党の差を出す。
野党の方が与党のほうが、一次情報の割合が高い傾向にある。


\textbf{Distribution of Dominance.}
まず、ペルソナ別 + alignment別をマージしたチャートを出す。
次に野党と与党の差を出す。

ChatGPTはGemini,ClaudeよりもPrimary informationの割合が高い傾向にある。
Claudeは全体的に、low barrierが回答に占める引用テキストの割合が高い傾向にある。


\textbf{Distribution of Citation Grounding.}
まず、ペルソナ別 + alignment別をマージしたチャートを出す。
次に野党と与党の差を出す。


全体、、階段になっている。

ChatGPTはPrimary informationの割合が高い傾向にある。また、distribution of citationよりも、一次情報の割合が高い傾向にある。primary informationの反映率は他のbarrierより強い傾向にある。
Dominanceも高いため、ChatGPTは一次情報を多く引用し、かつ引用テキストの割合も高い傾向にある。

一方、Claudeは反映率も、dominationに比例して高くなっている。
したがって攻撃面として、比較的発信コストが低いコンテンツが、回答の内容に大きな影響を与える可能性がある。